\documentclass[11pt,letterpaper]{article}

\usepackage[margin=0.85in]{geometry}
\usepackage[T1]{fontenc}
\usepackage[utf8]{inputenc}
\usepackage{lmodern}
\usepackage{microtype}
\usepackage{amsmath,amssymb,bm}
\usepackage{booktabs,tabularx,array}
\usepackage{graphicx}
\usepackage{xcolor}
\usepackage[most]{tcolorbox}
\usepackage{enumitem}
\usepackage{fancyhdr}
\usepackage{hyperref}
\usepackage{float}

\definecolor{navy}{HTML}{17324D}
\definecolor{teal}{HTML}{147D7E}
\definecolor{gold}{HTML}{D49A2A}
\definecolor{lightgold}{HTML}{FFF7E6}

\hypersetup{
  colorlinks=true,
  linkcolor=navy,
  urlcolor=teal,
  pdftitle={Informe del experimento de dimensión latente de un VAE}
}

\pagestyle{fancy}
\fancyhf{}
\lhead{Experimento de dimensión latente de un VAE}
\rhead{Procesamiento de Imágenes y Visión por Computadora}
\cfoot{\thepage}
\setlength{\headheight}{14pt}

\graphicspath{{outputs/latent_dimension_experiment/}}

\newcommand{\studentname}{Nombre del estudiante}
\newcommand{\studentid}{0000}
\newcommand{\coursesection}{Sección}
\newcommand{\submissiondate}{\today}

\newtcolorbox{responsebox}[1][]{
  colback=white,
  colframe=teal!70!black,
  boxrule=0.8pt,
  arc=2pt,
  title=Respuesta,
  fonttitle=\bfseries,
  breakable,
  #1
}

\newtcolorbox{evidencebox}[1][]{
  colback=lightgold,
  colframe=gold!80!black,
  boxrule=0.8pt,
  arc=2pt,
  title=Evidencia que debe incluir,
  fonttitle=\bfseries,
  breakable,
  #1
}

\newcommand{\answerarea}[1][2.1cm]{%
\begin{responsebox}
\textit{Escriba aquí su respuesta.}
\vspace{#1}
\end{responsebox}
}

\newcommand{\experimentfigure}[3]{%
\begin{figure}[H]
\centering
\IfFileExists{outputs/latent_dimension_experiment/#1}{%
  \includegraphics[width=#2\textwidth]{#1}%
}{%
  \fbox{\parbox[c][5.0cm][c]{0.88\textwidth}{\centering
  Ejecute el notebook para generar \texttt{\detokenize{#1}}.}}%
}
\caption{#3}
\end{figure}
}

\title{\textbf{Informe de laboratorio: selección de la dimensión latente de un autoencoder variacional}}
\author{}
\date{}

\begin{document}
\maketitle

\begin{center}
\renewcommand{\arraystretch}{1.35}
\begin{tabularx}{0.92\textwidth}{@{}>{\bfseries}lX>{\bfseries}lX@{}}
Nombre: & \studentname & Carné: & \studentid \\
Sección: & \coursesection & Fecha de entrega: & \submissiondate \\
\end{tabularx}
\end{center}

\section{Objetivo del laboratorio}

Se desea seleccionar un autoencoder variacional capaz de reconstruir dígitos
reconocibles, generar muestras plausibles a partir de la distribución normal
estándar, producir interpolaciones graduales y mantener una representación
compacta. Se comparan las dimensiones latentes
\[
d_z\in\{2,8,16,64\}.
\]
La recomendación final debe apoyarse en evidencia cuantitativa y cualitativa.
Ninguna métrica individual es suficiente para decidir.

\section{Hipótesis iniciales}

Complete esta sección antes de examinar los resultados.

\subsection*{Pregunta 1}
¿Qué dimensión latente espera que produzca la menor pérdida de reconstrucción?
Explique la relación esperada entre capacidad latente y reconstrucción.
\answerarea[2.5cm]

\subsection*{Pregunta 2}
¿Qué dimensión espera que produzca las muestras más plausibles para
$\mathbf z\sim\mathcal N(0,I)$? Explique por qué la calidad de muestreo puede no
seguir el mismo comportamiento que la reconstrucción.
\answerarea[2.5cm]

\subsection*{Pregunta 3}
Prediga si todas las dimensiones disponibles estarán activas en el modelo con
$d_z=64$. ¿Qué significa que una dimensión se encuentre inactiva?
\answerarea[2.3cm]

\subsection*{Pregunta 4}
Escriba una recomendación inicial. Indique qué evidencia podría hacerle cambiar
esa recomendación.
\answerarea[2.5cm]

\section{Integridad experimental}

\subsection*{Pregunta 5}
Registre la configuración y explique por qué deben mantenerse constantes todos
los parámetros distintos de la dimensión latente.

\begin{center}
\renewcommand{\arraystretch}{1.3}
\begin{tabular}{@{}ll@{}}
\toprule
Configuración & Valor \\
\midrule
Conjunto de datos & MNIST \\
Dimensiones latentes & $2,8,16,64$ \\
$\beta$ & \rule{3cm}{0.4pt} \\
Épocas & \rule{3cm}{0.4pt} \\
Tamaño de lote & \rule{3cm}{0.4pt} \\
Tasa de aprendizaje & \rule{3cm}{0.4pt} \\
Semilla aleatoria & \rule{3cm}{0.4pt} \\
Umbral de actividad & \rule{3cm}{0.4pt} \\
Dispositivo utilizado & \rule{3cm}{0.4pt} \\
\bottomrule
\end{tabular}
\end{center}
\answerarea[2.3cm]

\section{Resultados cuantitativos}

\IfFileExists{outputs/latent_dimension_experiment/summary_table.tex}{%
\begin{table}[H]
\centering
\resizebox{\textwidth}{!}{\begin{tabular}{rrrrrrrr}
\toprule
Latent dim & Parameters & Test negative ELBO & Test reconstruction & Test KL & Active dimensions & Active fraction & Training seconds \\
\midrule
2 & 88645 & 159.342 & 153.772 & 5.570 & 2 & 1.000 & 284.881 \\
8 & 145105 & 115.518 & 99.339 & 16.178 & 8 & 1.000 & 283.950 \\
16 & 220385 & 105.613 & 80.834 & 24.779 & 16 & 1.000 & 399.093 \\
64 & 672065 & 104.853 & 75.783 & 29.070 & 42 & 0.656 & 824.937 \\
\bottomrule
\end{tabular}
}
\caption{Comparación final de las dimensiones latentes candidatas.}
\end{table}
}{%
\begin{center}
\fbox{\parbox[c][3.2cm][c]{0.88\textwidth}{\centering
Ejecute el notebook para generar \texttt{summary\_table.tex}.}}
\end{center}
}

\subsection*{Pregunta 6}
Ordene los cuatro modelos según su pérdida de reconstrucción en prueba. ¿Las
mejoras entre dimensiones consecutivas son sustanciales o marginales? Incluya
los valores medidos.
\answerarea[2.8cm]

\subsection*{Pregunta 7}
¿Cómo cambian la divergencia KL total, el número de dimensiones activas y la
fracción activa al aumentar $d_z$? Explique qué revela cada medida sobre la
capacidad latente efectiva.
\answerarea[3.0cm]

\subsection*{Pregunta 8}
Compare el número de parámetros y el tiempo de entrenamiento. ¿Aumentar $d_z$
modifica sustancialmente el tamaño total y el costo computacional del modelo?
\answerarea[2.8cm]

\section{Comportamiento durante el entrenamiento}

\experimentfigure{training_curves.png}{0.98}{Comportamiento de entrenamiento y prueba para cada dimensión latente.}

\subsection*{Pregunta 9}
¿Todos los modelos alcanzaron una región razonablemente estable en la última
época? Identifique cualquier comparación afectada por convergencia incompleta.
\answerarea[2.5cm]

\subsection*{Pregunta 10}
Describa la relación entre la pérdida de reconstrucción, la divergencia KL y el
ELBO negativo. ¿Qué término explica las principales diferencias entre modelos?
\answerarea[2.8cm]

\section{Análisis de reconstrucción}

\experimentfigure{reconstructions.png}{0.98}{Reconstrucciones de las mismas imágenes de prueba mediante las medias posteriores.}

\subsection*{Pregunta 11}
Identifique al menos dos imágenes en las que aumentar $d_z$ produzca una mejora
visual significativa. Describa cambios en trazos, forma, ambigüedad o detalles.
\answerarea[3.0cm]

\subsection*{Pregunta 12}
¿A partir de qué dimensión las dimensiones adicionales muestran rendimientos
decrecientes? Utilice evidencia visual y la pérdida de reconstrucción.
\answerarea[2.8cm]

\section{Análisis de muestras del prior}

\experimentfigure{prior_samples.png}{0.98}{Muestras generadas a partir de $\mathbf z\sim\mathcal N(0,I)$.}

\subsection*{Pregunta 13}
Compare el realismo y la diversidad de las muestras. ¿Qué modelo produce la
mayor proporción de dígitos reconocibles? ¿Cuál muestra mayor variedad visible?
\answerarea[3.0cm]

\subsection*{Pregunta 14}
¿El modelo con la mejor reconstrucción también produce las mejores muestras del
prior? Explique por qué ambos resultados pueden diferir en un VAE.
\begin{evidencebox}
Considere el término de reconstrucción y la necesidad de que los posteriores
codificados sean compatibles con el prior utilizado para generar.
\end{evidencebox}
\answerarea[3.0cm]

\section{Análisis de interpolación}

\experimentfigure{interpolations.png}{0.98}{Interpolación lineal entre las medias posteriores del par seleccionado.}

\subsection*{Pregunta 15}
¿Qué modelo produce la transición más gradual y plausible? Describa dónde cambia
la identidad y si los trazos intermedios permanecen coherentes.
\answerarea[3.0cm]

\subsection*{Pregunta 16}
¿Una interpolación suave demuestra que cada dimensión latente tiene un significado
interpretable? Explique la limitación de esa conclusión.
\answerarea[2.6cm]

\section{Organización del espacio latente}

\experimentfigure{latent_space.png}{0.86}{Medias posteriores: representación directa para $d_z=2$ y PCA para espacios mayores.}

\subsection*{Pregunta 17}
¿Qué clases forman regiones distinguibles? Identifique clases que se superponen
y explique si su similitud visual puede justificarlo.
\answerarea[3.0cm]

\subsection*{Pregunta 18}
¿Por qué la gráfica directa de $d_z=2$ debe interpretarse de manera diferente a
las proyecciones bidimensionales obtenidas mediante PCA?
\answerarea[2.8cm]

\section{Análisis de dimensiones activas}

\experimentfigure{active_dimensions.png}{0.88}{Contribución KL promedio de cada dimensión; la línea discontinua indica el umbral de actividad.}

\subsection*{Pregunta 19}
¿Cada modelo utiliza todas las dimensiones disponibles? Compare el número y la
proporción de dimensiones activas.
\answerarea[2.7cm]

\subsection*{Pregunta 20}
Si el modelo de 64 dimensiones deja muchas dimensiones inactivas, ¿qué implica
para su capacidad nominal y efectiva? ¿Debe considerarse automáticamente un
fracaso de entrenamiento?
\answerarea[3.2cm]

\section{Recomendación final}

\subsection*{Pregunta 21: dimensión seleccionada}
\begin{center}
\Large Recomiendo $d_z=\boxed{\phantom{000}}$ para el escenario planteado.
\end{center}

\subsection*{Pregunta 22: justificación basada en evidencia}
Defienda su decisión utilizando:
\begin{itemize}[leftmargin=*]
  \item al menos tres valores cuantitativos;
  \item evidencia de reconstrucción;
  \item realismo y diversidad de las muestras del prior;
  \item comportamiento de las interpolaciones;
  \item dimensiones latentes activas; y
  \item compacidad o costo computacional.
\end{itemize}
\answerarea[7.0cm]

\subsection*{Pregunta 23: debilidad del modelo seleccionado}
Identifique la principal debilidad del modelo recomendado. Explique por qué la
recomendación sigue siendo razonable.
\answerarea[3.2cm]

\subsection*{Pregunta 24: escenario alternativo}
Describa un uso realista en el que otra dimensión sería preferible. Seleccione
la dimensión alternativa y justifique su elección.
\answerarea[3.5cm]

\section{Limitaciones y experimentos posteriores}

\subsection*{Pregunta 25}
Identifique al menos tres limitaciones del experimento. Considere duración del
entrenamiento, variación aleatoria, subjetividad visual, complejidad de MNIST,
arquitectura y umbral de actividad.
\answerarea[4.0cm]

\subsection*{Pregunta 26}
Proponga un experimento posterior que fortalezca la recomendación. Indique la
variable que cambiaría, las cantidades que mediría y qué resultado apoyaría o
contradiría su decisión.
\answerarea[4.0cm]

\section{Conclusión}

Escriba una conclusión de 150--250 palabras. Indique el modelo seleccionado,
resuma la evidencia decisiva y describa el principal compromiso observado.
\answerarea[4.0cm]

\section*{Lista de verificación}
\begin{itemize}[label=$\square$,leftmargin=*]
  \item Se respondieron las 26 preguntas.
  \item La recomendación final selecciona una dimensión latente.
  \item Las afirmaciones cuantitativas incluyen valores medidos.
  \item Las afirmaciones cualitativas hacen referencia a figuras o ejemplos.
  \item Se diferencia la gráfica directa de las proyecciones mediante PCA.
  \item Se analiza una limitación y un escenario alternativo.
  \item El informe completo compila sin errores.
\end{itemize}

\end{document}
